\section{小题2：SysY 等价 LLVM IR 与 RISC-V 程序实验}
小题2由两名成员分别完成 LLVM IR 与目标汇编两个任务。两部分共享同一份 SysY 示例程序和同一组测试数据，使两种实现围绕统一语义基准进行比较。SysY 程序由两人共同确定；成员B负责 LLVM IR，成员A负责 RISC-V 汇编。

\subsection{共同 SysY 示例程序设计}
SysY 是由 C 语言子集扩展而来的教学语言。本实验设计的程序读入 $n$ 个整数，计算其中所有偶数之和，并统计正数个数。程序包含数组、赋值、算术与取模运算、关系判断、\texttt{if} 条件分支、\texttt{while} 循环、函数、数组参数和函数调用等语言特性；输入输出通过 SysY 运行时库中的 \texttt{getint}、\texttt{putint} 和 \texttt{putch} 完成。为保证数组访问范围有效，测试时取 $0\le n\le 10$。

\begin{lstlisting}[language=C,caption={SysY 示例程序 main.sy},label={lst:sysy-main}]
//求偶数和
int sum_even(int a[],int n)
{
    int i;
    int sum;
    i=0;
    sum=0;
    while(i<n)
    {
        if(a[i]%2==0)
            sum=sum+a[i];
        i=i+1;
    }
    return sum;
}

//统计正数个数
int count_pos(int a[],int n)
{
    int i;
    int cnt;
    i=0;
    cnt=0;
    while(i<n)
    {
        if(a[i]>0)
            cnt=cnt+1;
        i=i+1;
    }
    return cnt;
}

int main()
{
    int a[10];
    int n;
    int i;
    int s;
    int c;

    //输入
    n=getint();
    i=0;
    while(i<n)
    {
        a[i]=getint();
        i=i+1;
    }

    //计算
    s=sum_even(a,n);
    c=count_pos(a,n);

    //输出
    putint(s);
    putch(10);
    putint(c);
    putch(10);

    return 0;
}
\end{lstlisting}

为便于两种实现统一验证，固定使用三组测试数据。第1组同时包含正数、负数、奇数和偶数；第2组验证负偶数参与求和；第3组全部为正奇数，用于验证偶数和为 0 的情况。两名成员的 LLVM IR 与 RISC-V 实现均应使用这三组输入并得到相同输出。

\begin{table}[H]
\centering
\caption{小题2统一测试数据}
\begin{tabularx}{0.92\linewidth}{c X c c}
\toprule
组别 & 输入（第一项为 $n$） & 偶数和 & 正数个数 \\
\midrule
1 & 5；1 2 -3 4 6 & 12 & 4 \\
2 & 4；-2 3 8 -5 & 6 & 2 \\
3 & 5；1 3 5 7 9 & 0 & 5 \\
\bottomrule
\end{tabularx}
\end{table}
